\chapter{INTRODUCTION}

\section{Project Overview}

Public transportation is the backbone of mobility in modern cities, and the reliability of the information provided to passengers directly shapes their travel experience. One of the most frequently asked questions at any bus stop is simply: \textit{``When will my bus arrive?''} Accurate answers to this question reduce perceived waiting time, increase passenger satisfaction, and encourage the use of public transit over private vehicles. Consequently, Bus Arrival Time Prediction (BATP) has become a core component of Advanced Traveler Information Systems (ATIS) in intelligent transportation.

This project develops and evaluates a context-aware machine learning system that predicts inter-stop bus travel times for the İzmir ESHOT public transportation network. The system combines real-time Global Positioning System (GPS) data, static General Transit Feed Specification (GTFS) schedule data, and contextual signals such as weather and traffic to estimate how long a bus takes to move from one stop to the next. These segment-level predictions can be accumulated along a route to estimate the arrival time at any downstream stop.

The work is anchored on a real, self-collected dataset. A custom data collector continuously queried the İzmir open transit Application Programming Interface (API) over a window spanning April to June 2026, recording several million GPS observations across three bus routes. From these observations, clean stop-to-stop travel-time segments were extracted and enriched with engineered features. A range of models was then trained and compared, from simple baselines (a naïve schedule-based predictor and a historical average) to classical machine learning models (Random Forest and XGBoost) and a deep learning model (a dual-input Long Short-Term Memory network, LSTM). The methodology is reproducible end to end and every reported result is fixed by a random seed.

The reference point for this study is the work of Kaya and Kalay \cite{kaya2025}, who proposed context-aware deep learning models for spatio-temporal bus arrival time forecasting and reported a Mean Absolute Error (MAE) of 2.97 minutes. Our project reproduces a comparable pipeline on independently collected İzmir data and, rather than only chasing a lower error number, places strong emphasis on \textit{methodological honesty}: each design choice is validated through controlled A/B experiments, and several widely-assumed ``improvements'' are shown empirically not to help.

\section{Motivation and Problem Statement}

\subsection{Importance of Accurate Arrival Prediction}

Real-time arrival information is one of the highest-value features a transit agency can offer. Studies on traveler behavior consistently show that uncertainty, rather than the absolute waiting time, is the dominant source of passenger frustration. When passengers can see a trustworthy estimate of their bus's arrival, the same physical wait feels shorter and more tolerable. For the operator, accurate predictions enable better fleet monitoring, schedule adherence analysis, and operational planning.

İzmir, the third-largest city in Türkiye, operates an extensive bus network through ESHOT. Although the municipality exposes real-time bus positions through a public API, it does not provide a GTFS-Realtime feed, and there is no readily available, openly documented arrival-prediction model tuned to local conditions. This gap is the practical motivation for the project.

\subsection{Problem Statement}

The central problem addressed in this project is the prediction of the travel time of a bus between two consecutive stops, given the bus's recent trajectory and the prevailing temporal, spatial, and environmental context. Formally, for a bus traveling from stop $i$ to stop $i+1$, the goal is to predict the travel time $t_{i \rightarrow i+1}$ using features available \emph{before} the bus reaches stop $i+1$, so that the prediction is causally valid and free of information leakage.

Three properties make this problem challenging. First, inter-stop travel times in dense urban routes are short (on the order of one minute) and highly variable, dominated by traffic signals, passenger boarding, and congestion. Second, the only freely available position signal is polled at a coarse 30-second interval, which imposes a quantization floor on how precisely the ground-truth travel time can even be measured. Third, the contextual features that are theoretically helpful (such as weather) may behave as noise when the collected data does not contain enough variation in those conditions.

\section{Objectives and Scope}

\subsection{Primary Objectives}

The development of this project is guided by the following objectives:

\begin{enumerate}
    \item \textbf{Build an end-to-end, reproducible BATP pipeline:} from real-time data ingestion and trip extraction, through feature engineering, to model training and evaluation.
    \item \textbf{Exploit schedule information as a feature:} use GTFS scheduled travel times and the deviation from schedule as explicit model inputs, which is a central design hypothesis of this work.
    \item \textbf{Compare classical and deep learning models fairly:} benchmark baselines, Random Forest, XGBoost, and an LSTM on the \emph{same} test set, and quantify any differences with statistical significance tests.
    \item \textbf{Validate every design decision with controlled experiments:} adopt an A/B methodology in which each proposed change is kept only if it measurably improves the result.
    \item \textbf{Demonstrate generalization:} confirm that the same method works on more than one route, rather than overfitting to a single line.
\end{enumerate}

\subsection{Scope}

The scope of this project encompasses the following areas:

\textbf{Data Collection:}
\begin{itemize}
    \item Real-time GPS polling of ESHOT buses every 30 seconds via the İzmir open transit API.
    \item Hourly weather acquisition and periodic traffic-flow sampling as contextual signals.
    \item Persistent storage of all observations in a relational database for offline processing.
\end{itemize}

\textbf{Feature Engineering:}
\begin{itemize}
    \item Temporal features (hour, day type, cyclic encodings).
    \item Spatial features (inter-stop distance, progress along the route).
    \item Schedule features derived from GTFS (scheduled travel time, deviation).
    \item Historical and lag features, and dwell-time features derived from raw GPS.
\end{itemize}

\textbf{Modeling and Evaluation:}
\begin{itemize}
    \item Five baseline and improved models plus a dual-input LSTM.
    \item Evaluation against the reference paper, condition-based analysis, and statistical significance testing.
    \item Generalization across three İzmir routes (502, 268, 565).
\end{itemize}

\subsection{Limitations}

The following limitations are explicitly acknowledged and are revisited in the conclusion:

\begin{itemize}
    \item \textbf{Measurement floor:} the 30-second public GPS polling interval limits the precision of the ground-truth travel times; finer labels would require denser positioning data that the public API does not provide.
    \item \textbf{Weather coverage:} the collected data is dominated by clear and cloudy conditions, with only a small number of rainy segments and no snow; although rainy segments do show higher error, the sample is too small to train a reliable weather-aware model.
    \item \textbf{No live demonstrator:} the project focuses on the predictive pipeline; a passenger-facing live dashboard is left as future work.
\end{itemize}

\section{Contributions}

The main contributions of this project are:

\begin{enumerate}
    \item A fully reproducible, real-time data collection and prediction pipeline for the İzmir ESHOT network, built on independently gathered data.
    \item Empirical evidence that GTFS schedule-derived features and GPS-derived dwell-time features are among the most informative inputs for inter-stop travel-time prediction.
    \item A rigorous, A/B-driven evaluation that yields several \emph{methodological findings}: classical XGBoost is statistically equivalent to a deep LSTM, the achievable error is bounded by a measurement (polling) floor, and additional model complexity does not necessarily improve accuracy.
    \item A fair, same-test-set statistical comparison of classical and deep models, correcting an earlier ``deep learning is best'' conclusion that turned out to be a test-set artifact.
    \item A demonstration of generalization across three routes with different traffic profiles.
\end{enumerate}

\section{Report Organization}

The remainder of this report is organized as follows. Chapter~2 reviews background concepts and related work, including the reference study. Chapter~3 describes the system design and methodology, covering data collection, feature engineering, and model architectures. Chapter~4 details the implementation, including the improvement program and the A/B evaluation methodology. Chapter~5 presents the experimental results, the statistical comparison of models, the generalization study, and the comparison with the literature. Chapter~6 concludes the report with a summary of findings, limitations, and directions for future work.
